\begin{circuitikz}[
    american,
    voltage dir=RP,
    >=latex,
    %>={Latex[length=5mm, width=2mm]},
    alt={Ladder logic diagram titled Toggle Logic - Second Scan featuring two rungs bounded by vertical power rails with green highlighting indicating active power flow. On the top rung, power flow highlighting extends only through the examine if closed contact labeled AutoPB, stopping before the OneShot instruction and leaving output coil AutoSS unhighlighted. On the bottom rung, power flow bypasses the top branch and routes down through the lower seal-in branch containing an examine if open contact labeled AutoSS in series with an examine if closed contact labeled Auto, before continuing straight across to energize the output coil labeled Auto.}
    ]

    % Define the number of rungs in the diagram
    \def\numRungs{2}

    % Calculate the length of the vertical power rails dynamically based on number of rungs
    \pgfmathsetmacro{\railLength}{-2 * \numRungs}

    %======================================================================================================
    % Component Grid
    % Spacing = 3x, 2y
    %======================================================================================================
    \foreach \i in {0,...,5} {
        \foreach \j in {0,...,20} { % Change to 0,...,19 if you need 20 individual Y points down to -40
        
        % Calculate spatial positions
        \pgfmathsetmacro{\px}{\i * 3}
        \pgfmathsetmacro{\py}{-1 - (\j * 2)}
        
        % Name coordinate using loop indices: (C-column-row) e.g., (C-0-0), (C-5-10)
        \coordinate (C-\i-\j) at (\px, \py);
        
        % OPTIONAL: Uncomment the 2 lines below to display grid labels while building
        % \fill[red] (C-\i-\j) circle (1.5pt);
        %\node[anchor=south west, scale=0.5, gray] at (C-\i-\j) {(\i,\j)};
        }
    }

    %======================================================================================================
    % Foreground Layer
    % Only contains the vertical power rails
    % Length of the rails is automatic, calculated based on the number of rungs defined above
    %======================================================================================================
    \begin{pgfonlayer}{ft}

        %======================================================================================================
        % Scope sets options for both lines - BakerBlack, very thick
        %======================================================================================================
        \begin{scope}[
            draw=BakerBlack,
            text=BakerBlack,
            color=BakerBlack,
            font=\small,
            ultra thick
        ]
        
            \draw (0,0) -- (0,\railLength);

            \draw (15,0) -- (15,\railLength);

        \end{scope}
    \end{pgfonlayer}

    %======================================================================================================
    % Main Layer
    % All actual logic goes here
    %======================================================================================================
    \begin{pgfonlayer}{main}
        
        %======================================================================================================
        % Scope sets options for both lines - BakerBlack for draw and text
        %======================================================================================================
        \begin{scope}[
            draw=BakerBlack,
            text=BakerBlack,
            color=BakerBlack,
            font=\small
        ]

            %=================================================================
            % Rung 1
            %=================================================================
            \draw (C-0-0) to[C, l={AutoPB}] (C-1-0) pic (ONS) {ONS};
            \draw (ONS-OUT) -- (C-4-0) to[esource, l={AutoSS}] (C-5-0);
            \draw (C-0-1) to[C, l={AutoSS}] (C-1-1) to[vC, l={Auto}] (C-2-1) -- (C-4-1) to[esource, l={Auto}] (C-5-1);
            \draw ($(C-0-1)+(0.25,0)$) coordinate(h1) -- ++(0,-2) coordinate(h2) to[vC, l={AutoSS}] (C-1-2) to[C, l={Auto}] ($(C-2-2)-(0.25,0)$) coordinate(h3) -- ++(0,2) coordinate(h4);

            \node[anchor=center] at (ONS-INS) {ONS};
            \node[anchor=south] at (ONS-VAR) {OneShot};

        \end{scope}
    \end{pgfonlayer}

    %=================================================================
    % Background layer
    %=================================================================
    \begin{pgfonlayer}{bg}
        
        %==================================================================
        % Highlight Scope
        % This layer exists for highlighting true logic, the scope sets the
        % highlight options without having to repeat them every \draw
        % draw=UpNorthGreen, line width=10pt, opacity=0.4
        %==================================================================
        \begin{scope}[
            draw=UpNorthGreen,
            line width=10pt,
            opacity=0.4
        ]

            %==================================================================
            % Left vertical rung - leave this alone, it should ALWAYS be green
            % It automatically draws itself to the length defined at the top
            %==================================================================
            \draw (0,0) -- (0,\railLength);
        
            \draw (C-0-0) -- (C-1-0);

            \draw (C-0-1) -- (h1) -- (h2) -- (h3) -- (h4) -- (C-5-1);

        \end{scope}
    \end{pgfonlayer}

\end{circuitikz}